% Generated draft: requires mathematical and source review.
\documentclass[11pt]{article}
\usepackage[margin=25mm]{geometry}
\usepackage{fontspec}
\setmainfont{Latin Modern Roman}
\usepackage{amsmath,amssymb,mathtools,unicode-math}
\setmathfont{Latin Modern Math}
\usepackage{enumitem,needspace,xurl,hyperref}
\setlength{\parindent}{0pt}
\setlength{\parskip}{5pt}
\setlength{\emergencystretch}{3em}
\newcommand{\N}{\mathbb{N}}
\newcommand{\Z}{\mathbb{Z}}
\newcommand{\Q}{\mathbb{Q}}
\newcommand{\R}{\mathbb{R}}
\newcommand{\C}{\mathbb{C}}
\newcommand{\F}{\mathbb{F}}
\newcommand{\A}{\mathbb{A}}
\newcommand{\PP}{\mathbb{P}}
\newcommand{\E}{\mathbb{E}}
\newcommand{\eps}{\varepsilon}
\newcommand{\dd}{\,\mathrm{d}}
\begin{document}
\section*{UnsolvedMath ID: 30005885}
\textbf{Source problem code:} catalogue entry
\subsection{Definitions and mathematical statement}
\begin{enumerate}
\item For a field \(K\) and a positive integer \(d\), let \(M_d(K)\) denote the \(d\times d\) matrices over \(K\), and let \(\mathrm{GL}_d(K)\) denote the invertible matrices. Define the normalized rank distance by
\[
 d_{\mathrm{rk}}(S,T)=\frac{\operatorname{rank}_K(S-T)}{d}.
\]
The source asks whether almost commuting invertible matrices admit nearby commuting invertible matrices in this metric. The scalar field, the uniformity across fields, and the permitted dimensions of the approximating matrices are not specified. With a fixed field \(K\) and approximants of the same dimension, the usual uniform stability formulation would ask whether
\[
\begin{gathered}
\forall\varepsilon>0\ \exists\delta>0\ \forall d\geq1\ \forall A,B\in\mathrm{GL}_d(K),\\
 d_{\mathrm{rk}}(AB,BA)<\delta
 \ \Longrightarrow\ 
 \exists C,D\in\mathrm{GL}_d(K):\\
 CD=DC,\qquad
 d_{\mathrm{rk}}(A,C)<\varepsilon,\qquad
 d_{\mathrm{rk}}(B,D)<\varepsilon.
\end{gathered}
\]
Here \(\delta\) is independent of \(d,A,B\); its possible dependence on \(K\) remains unresolved. This display is an explicitly conditional interpretation of the source question, requiring confirmation of the field and approximation conventions.

\item For a finite group \(G\), define its exponent by
\[
\exp(G)=\min\{e\in\Z_{\geq1}:g^e=1_G\text{ for every }g\in G\}.
\]
Is the following assertion true? For every integer \(k\geq2\) and every prime integer \(n\), there exists an integer \(N=N(k,n)\geq1\) such that, for every finite group \(G\) and every generating set \(A=\{a_1,\ldots,a_k\}\) of \(G\),
\[
\left[
 \forall(x_1,\ldots,x_N)\in(A\cup\{1_G\})^N,
 \quad (x_1\cdots x_N)^n=1_G
\right]
\quad\Longrightarrow\quad
\exp(G)\leq n?
\]
The integer \(N\) must be independent of \(G\) and its chosen generators. Since identity elements may be inserted, the hypothesis also covers all positive words of length at most \(N\) in the specified generators; inverses are not included unless they already belong to \(A\cup\{1_G\}\). The inequality on the exponent is retained as supplied. The original extract says ``any (prime) number \(n\)'', whereas the cleaned registry statement explicitly restricts \(n\) to primes; the intended scope beyond primes remains unclear.
\end{enumerate}
The supplied dated curation reports partial results for unitary or self-adjoint matrices and for finite solvable groups, and describes the arbitrary-invertible-matrix and general finite-group assertions as unresolved. These are supplied literature assessments, not independently verified mathematical conclusions here.
\subsection{Short English statement}
The first question asks whether two invertible matrices whose products almost agree in normalized rank distance can be approximated by commuting invertible matrices, uniformly in dimension. Its field and approximation conventions need clarification. The second asks whether, for fixed generator count \(k\) and prime \(n\), a bounded test on positive words forces every finite group passing the test to have exponent at most \(n\).
\subsection{Sources}
\begin{itemize}
\item Supplied local UnsolvedMath registry record, permanent ID \(30005885\), title \emph{Rank-Metric Stability and Finite-Group Exponent Problems}, exact source problem code \texttt{catalogue entry}; collection: \emph{Oberwolfach Reports Open Problems}; proposed year: 2024. The registry supplies no proposer in its structured attribution field.
\item Supplied report citation: \emph{Mini-Workshop: Growth and Expansion in Groups} (2024), Oberwolfach Reports, pp. 1005--1040, DOI 10.4171/owr/2024/17. \url{https://doi.org/10.4171/owr/2024/17}.
\item Supplied literature source, labeled ``rank-metric stability'': \url{https://ems.press/journals/ggd/articles/2152717}.
\item Supplied literature source, labeled ``OWR report'': \url{https://ems.press/content/serial-article-files/49475}.
\item Literature assessment date: \texttt{2026-08-22}. This is supplied dated curation, not independently verified. It labels the record partially solved, reports rank-metric stability for unitary or self-adjoint matrices and the exponent implication for finite solvable groups, and describes the arbitrary invertible and general finite-group cases as unresolved. The supplied background explicitly requests independent verification of time-sensitive status and cited claims.
\item The original extract contains a damaged attribution to Ol'shanskii adjacent to the finite-group question. It subsequently contains ``Questions by J. Bajpai'' and material on thin groups and free subgroups; that placement does not establish Bajpai's authorship of either question formalized here.
\item The registry labels its statement \texttt{corrected\_verified} and states that the official report was checked and the question isolated from surrounding exposition. This is supplied verification metadata; no source visit or independent verification is claimed in this draft.
\end{itemize}
\end{document}
